% Work-MAT | 2. rocnik | pripravovana hodina 35
% Návaznost: uživatelské připravy.zip (hodiny 01–20), autorská učebnice na GitHubu.
% Realisticky.cz: didaktická inspirace v posloupnosti a typech úloh; zadání jsou vlastní.
\documentclass[11pt,a4paper]{article}
\input{../../sablony/pripravy/priprava.sty}
\nastavPripravu{2. ročník}{Logaritmická funkce}{35}{Graf logaritmické funkce}
\begin{document}
\blok{Motivace a návaznost}
Přelož graf $y=2^x$ osovou souměrností podle $y=x$ a připomeň převrácení definičního oboru s oborem hodnot. Potom žáci kreslí $\log_2x$ a $\log_3x$.
\poznamka{Orientační nabídka pro 45--55 minut: krátké zopakování a motivace, nový výklad, tři společné vzorové úlohy, samostatná práce. Vyber podle tempa třídy; zbývající zadání slouží jako rezerva.}
\blok{Výklad}
Logaritmická funkce $f(x)=\log_a x$ má $D=(0;\infty)$, $H=\mathbb R$, prochází bodem $[1;0]$ a má svislou asymptotu $x=0$.
\blok{Společné příklady}
\begin{prikladReseniBox}[title={Příklad 1}]
Vyřeš, resp. urči: $f(x)=\log_2(x-2)$.
\tcblower
\[\begin{aligned}x-2>0\;\Longrightarrow\;D=(2;\infty) \\ \text{Svislá asymptota: }x=2\end{aligned}\]
\end{prikladReseniBox}
\begin{prikladReseniBox}[title={Příklad 2}]
Vyřeš, resp. urči: $f(x)=\log_3(x+1)$.
\tcblower
\[\begin{aligned}D=(-1;\infty) \\ \text{Asymptota: }x=-1\end{aligned}\]
\end{prikladReseniBox}
\poznamka{V grafických úlohách nech určovat osy včetně asymptot, úsečku mezi body $[1;0]$ a $[a;1]$, případné posuny.}
\blok{Samostatné procvičování}
\begin{zadaniReseniBox}[title={Samostatná práce 1}]
\begin{enumerate}[label=\alph*)]
\item $f(x)=\log_{2}(x)$.
\item $f(x)=\log_{3}(x)$.
\item $f(x)=\log_{2}(x+1)$.
\item $f(x)=\log_{2}(x-2)$.
\item $f(x)=\log_{3}(x+4)$.
\item $f(x)=\log_{5}(x)$.
\end{enumerate}
\tcblower
\begin{enumerate}[label=\alph*)]
\item $D(f)=(0;\infty),\;\text{asymptota }x=0$.
\item $D(f)=(0;\infty),\;\text{asymptota }x=0$.
\item $D(f)=(-1;\infty),\;\text{asymptota }x=-1$.
\item $D(f)=(2;\infty),\;\text{asymptota }x=2$.
\item $D(f)=(-4;\infty),\;\text{asymptota }x=-4$.
\item $D(f)=(0;\infty),\;\text{asymptota }x=0$.
\end{enumerate}
\end{zadaniReseniBox}
\blok{Dobrovolné rozšíření}
\begin{zadaniReseniBox}[title={Rozšíření 1}]
\begin{enumerate}[label=\alph*)]
\item $Vytvoř vlastní úlohu tohoto typu a se spolužákem porovnej podmínky a výsledek.$.
\end{enumerate}
\tcblower
\begin{enumerate}[label=\alph*)]
\item Zkontroluj definice a argumenty..
\end{enumerate}
\end{zadaniReseniBox}
\blok{Úkol na doma}
U každého logaritmu vypiš podmínky argumentu a proveď zkoušku. U rovnic vyznač případně nepřípustné kořeny, u nerovnic odpověz intervalem.
\begin{zadaniReseniBox}[title={Domácí úkol}]
\begin{enumerate}[label=\alph*)]
\item $f(x)=\log_{2}(x+3)$.
\item $f(x)=\log_{3}(x-1)$.
\item $f(x)=\log_{5}(x+2)$.
\item $f(x)=\log_{2}(x-4)$.
\item $f(x)=\log_{4}(x+1)$.
\item $f(x)=\log_{3}(x+2)$.
\item $f(x)=\log_{2}(x+5)$.
\item $f(x)=\log_{10}(x-2)$.
\end{enumerate}
\tcblower
\begin{enumerate}[label=\alph*)]
\item $D(f)=(-3;\infty),\;\text{asymptota }x=-3$.
\item $D(f)=(1;\infty),\;\text{asymptota }x=1$.
\item $D(f)=(-2;\infty),\;\text{asymptota }x=-2$.
\item $D(f)=(4;\infty),\;\text{asymptota }x=4$.
\item $D(f)=(-1;\infty),\;\text{asymptota }x=-1$.
\item $D(f)=(-2;\infty),\;\text{asymptota }x=-2$.
\item $D(f)=(-5;\infty),\;\text{asymptota }x=-5$.
\item $D(f)=(2;\infty),\;\text{asymptota }x=2$.
\end{enumerate}
\end{zadaniReseniBox}
\poznamka{Domácí úkol obsahuje 8 položek; předpokládaná domácí práce přibližně 15--30 minut. Po větších blocích jsou vědomě promíchané dříve probrané typy.}
\end{document}
